% ECONOMICS_PROBLEM: {"id": "C90-643", "title": "Experimenter demand and elicitation", "jel_code": "C90", "source_id": "OP-0307"}
% !TeX program = xelatex
\documentclass[11pt,a4paper]{article}
\usepackage[margin=23mm]{geometry}
\usepackage{fontspec}
\setmainfont{Latin Modern Roman}
\usepackage{amsmath,amssymb,mathtools}
\usepackage{xcolor,enumitem,booktabs,longtable,array,xurl,hyperref}
\definecolor{muted}{HTML}{53636C}
\hypersetup{colorlinks=true,urlcolor=blue,linkcolor=blue}
\setlength{\parindent}{0pt}
\setlength{\parskip}{5pt}
\newcommand{\R}{\mathbb R}
\newcommand{\E}{\mathbb E}
\newcommand{\N}{\mathbb N}
\newcommand{\ind}{\mathbf 1}
\newcommand{\Var}{\operatorname{Var}}
\newcommand{\Cov}{\operatorname{Cov}}
\newcommand{\supp}{\operatorname{supp}}
\DeclareMathOperator*{\argmin}{arg\,min}
\DeclareMathOperator*{\argmax}{arg\,max}
\newcommand{\entrymeta}[3]{{\sffamily\small\color{muted}\textbf{\texttt{#1}}\quad|\quad #2\par #3\par}}
\newcommand{\Problemref}[1]{\texttt{#1}}
\newcommand{\JELref}[2]{\texttt{#2} (JEL \texttt{#1})}
\begin{document}
\section*{Experimenter demand and elicitation}
\textbf{Problem ID:} \texttt{C90-643}\quad \textbf{JEL:} \texttt{C90}\par
% BEGIN ECONOMICS STATEMENT
\label{problem:OP-0307}
{\sffamily\small\color{muted}Original subject group: Behavioral and experimental economics\par}
\label{definitions:problem:30007589}
\textbf{Audit classification:} Published research agenda. de Quidt–Haushofer–Roth §VI explicitly proposes mitigation and environmental/incentive studies. The monotone bracket is an existing method’s editorial specialization, not a newly open bound.
\subsection*{Definitions and precise research target}
\textbf{Model and research target.} Let \(D\in\{0,1\}\) be an experimentally assigned substantive treatment and \(q\in[-1,1]\) a participant's belief about the action favored by the researcher, with \(q=0\) denoting a specified neutral-demand benchmark. Let \(Y_i(d,q)\) be the potential response of participant \(i\). Randomized demand cues \(Z\in\{-,0,+\}\) shift this belief while preserving the substantive task. The demand-free treatment estimand is
\[\tau^0=\mathbb E[Y_i(1,0)-Y_i(0,0)].\]
If negative and positive cues bracket neutral demand and \(Y_i(d,q)\) is monotone in \(q\), define \(\mu_d^-\) and \(\mu_d^+\) as mean outcomes under the negative and positive cues. Then the candidate bound is
\[\mu_1^- -\mu_0^+\leq\tau^0\leq\mu_1^+ -\mu_0^-.\]
Determine which cue and elicitation designs make these bracketing and exclusion restrictions credible across experimental environments, and which features reduce demand sensitivity without changing substantive preferences or information processing. Observed cue effects alone cannot verify that unobserved natural demand lies within the bracket. Measure cue comprehension, preregister effect-direction assumptions, and examine stronger cues, real incentives, anonymity, and survey versus laboratory contexts. Report partial identification when neutral behavior is not observed.

\emph{Definitions and maintained assumptions (editorial unless a source definition is named).} For each participant and substantive treatment \(d\), specify cue-induced beliefs \(q_i^-(d)\le0\le q_i^+(d)\) almost surely. Cue exclusion requires the observed response under cue \(z\) to equal \(Y_i(d,q_i^z(d))\), with no other utility, information or task change. Assume \(q\mapsto Y_i(d,q)\) is weakly increasing for every \(i,d\), or reverse orientations before using the bracket; an average cue shift or average monotonicity is not sufficient for the individual bracket. Randomize \((D,Z)\) with positive probability for every cell. Then \(\mu_d^-\le\mathbb E[Y_i(d,0)]\le\mu_d^+\), so the displayed interval follows by subtraction. This proof under maintained assumptions is elementary and already represented by the source’s bounding framework. Identifying the neutral effect requires observed neutral beliefs or stronger assumptions. Observed cue responses cannot prove the unobserved neutral belief is bracketed. Simultaneous sampling bounds are needed for the four means.
\subsection*{Variants and assumption changes}
\begin{enumerate}
\item \textbf{Monotone bound (source method, editorial specialization).} Use stronger positive/negative cues and randomized incentive levels, checking comprehension. Estimate valid confidence intervals under the individual bracketing and exclusion assumptions.
\item \textbf{Imperfect bracket (editorial).} Permit at most fraction \(\eta\) of participants to violate bracketing or monotonicity and bound \(Y\) in \([a,b]\). Derive sensitivity intervals widening the mean bounds by at most \(\eta(b-a)\) per treatment arm, and validate mitigation devices in new environments.
\end{enumerate}
\subsection*{References and source audit}
\begin{enumerate}
\item §VI pp.35–36 in April2018 manuscript. Supports the stated research area; exact displayed specification is editorial. \url{https://jondequidt.com/pdfs/deQuidt_Haushofer_Roth_demand_final.pdf}
\end{enumerate}
\begin{enumerate}\raggedright
\item\hypertarget{definitions:source:30007589:1}{} Jonathan de Quidt; Johannes Haushofer; Christopher Roth. 2018. \emph{Measuring and Bounding Experimenter Demand}. Author manuscript dated April 4, 2018, §VI Conclusion, printed pp. 35–36 (PDF pp. 35–36), particularly last paragraph..
\url{https://jondequidt.com/pdfs/deQuidt_Haushofer_Roth_demand_final.pdf}.
DOI: \url{https://doi.org/10.1257/aer.20171330}.
\end{enumerate}

\subsection*{Common conventions: Source mathematical conventions}

These conventions apply to each entry unless explicitly replaced.

\textbf{Models and identification.} A primitive model \(m\) specifies all probability laws, feasible choices, information, equilibrium and measurement rules used by its target. The admissible class \(\mathcal M\) is fixed independently of outcomes. If \(P_O(m)\) is its observable probability law and \(\tau(m)\) the target, its identified set at observable law \(P\) is
\[
 \mathcal I_\tau(P)=\{\tau(m):m\in\mathcal M,\ P_O(m)=P\}.
\]
Point identification means that this set is a singleton. An empty set signals incompatibility of data and assumptions. Coordinate bounds need not describe the joint set. Bounds are sharp only if every retained target is attainable or arbitrarily approachable by compatible models. Finite bounds require explicit restrictions on unobserved quantities; unrestricted confounding, hidden wealth, extrapolation or arbitrary equilibrium selection can yield uninformative sets.

\textbf{Causal interventions.} A potential outcome \(Y_i(p)\) is the outcome under a complete policy regime \(p\), including timing, financing, information and market spillovers. The target population and units are fixed before treatment. With interference, use the complete assignment vector or a justified exposure mapping; individual no-interference notation alone cannot identify an economy-wide intervention. A difference of mean potential outcomes does not require the cross-world joint distribution. Conditioning on post-treatment migration, survival, enrollment or employment changes the estimand and needs separate justification.

\textbf{Statistical statements.} An estimator, confidence set, forecast or test specifies a sampling law, model class and loss/coverage criterion. Uniform \((1-\alpha)\)-coverage means \(\inf_{m\in\mathcal M}\Pr_m\{\tau(m)\in C_n\}\geq1-\alpha\), or an explicitly asymptotic version. The whole parameter space trivially has coverage, so informativeness requires a stated width, risk or power objective. A forecasting improvement is relative to a named benchmark, information set and evaluation population.

\textbf{Equilibrium and optimization.} An equilibrium comprises feasible plans, optimality, beliefs where needed, budget constraints and all market/resource clearing. The primitives or a stated benchmark define each correspondence \(\mathcal E(p;m)\). Nonemptiness is assumed or proved, never inferred just from compact policy sets. A selected-equilibrium target uses a declared selection rule; a robust target quantifies over all permitted equilibria. Use suprema or infima unless attainment is established. An \(\varepsilon\)-optimal policy is feasible and within \(\varepsilon\) of the finite optimum value. Report an empty feasible set separately.

\textbf{Welfare and accounting.} Normative utility, interpersonal normalization, social weights, numeraire, discounting and the use of public receipts are primitives. Behavioral utility need not coincide with the welfare criterion. Household welfare includes ownership income and rents once; adding profits again to compensating variation generally double-counts them. A monetary equivalent is defined by an explicit transfer and price convention, with monotonicity and range conditions. Average effects, mechanism contributions, transfers and welfare losses are different targets. Infinite sums require convergence and dynamic feasible plans require terminal or no-Ponzi conditions.

\textbf{Source versions.} A working-paper date, revision date and journal publication year are distinguished. A DOI for a journal version is not automatically the DOI of an earlier manuscript. Locators refer to the stated version and printed pages where specified. A metadata-only check does not verify a claim about a full-text conclusion. Every entry's source subsection narrows the provenance claim accordingly.
% END ECONOMICS STATEMENT
\end{document}
